\section{Introduction}\label{sec:intro}

A linear Markov model fitted to the 5- and 20-day realized volatility of the S\&P 500 reports an entropy production of \EpWinObs per day. Fitted to surrogate data whose volatility is time reversible by construction, the same model reports \EpWinSurMed. Most of what the fitted model calls dissipation is made by the way the state was observed: two trailing averages of one signal, the shorter of which systematically leads the longer. Trailing windows of this kind are the standard preprocessing of early-warning indicators for critical transitions \cite{scheffer2009,kuehn2011,dakos2012,lenton2012}, and the currents fitted to windowed states are not idle: time irreversibility, measured by the difference between forward and backward cross-correlations, has been proposed as an early-warning signal \cite{kooloth2026,xu2023}, and forecasts of the committor type integrate a drift fitted to such a state \cite{xing2026b}; both inherit whatever the observer adds. The fitted model is misspecified---a windowed state is not Markov---but its drift and current are what such pipelines compute, so its irreversibility is the quantity to account for. We ask what trailing windows do to a stochastic transition, and how much of the irreversibility of a windowed state belongs to the system.

Two effects are at stake. Windows compress the state and destroy transition information that the latent process still carries, which limits what any early-warning statistic can detect and how skillful it can be \cite{boettiger2012,ditlevsen2010,duplessis2026}. And a coarse-graining that does not commute with time reversal can create apparent irreversibility, whereas one that commutes bounds the true entropy production from below \cite{martinez2019,hartich2024,bisker2024,esposito2012}; conversely, scalar observation of a linear system can hide its currents \cite{lucente2022prr}. Trailing windows do not commute with time reversal. The same construction is the belief rule of extrapolative investors in behavioral finance \cite{bgjs2015}, an analogy we use only for orientation.

Our results are four. (i) For linear filters of any stationary multichannel input, the lag-one antisymmetry---the discrete-time area-enclosing rate, which sets the current of a fitted linear Markov model---splits exactly into an observer term and a physical term, and time reversal of the data separates them without a model (Sec.~\ref{sec:theory}, Proposition~\ref{prop:decomp}). (ii) The entropy production of a fitted two-dimensional model factorizes into the squared antisymmetry and an innovation-geometry factor, which explains why reversible inputs can look strongly dissipative (Proposition~\ref{prop:factor}). (iii) The information lost to compression is a closed-form series for an Ornstein--Uhlenbeck input (Lemma~\ref{lem:mehler}) and is governed by two rates, local curvature and Kramers switching, in a double well (Sec.~\ref{sec:info}). (iv) In daily S\&P 500 data, time reversal, centered windows and sign randomization identify the observer term in the volatility windows and the leverage effect as the physical term; a stochastic-volatility model with leverage accounts for both, and the same model shows that the market's windows lose about half of the information that the latent volatility carries about its next excursion (Sec.~\ref{sec:market}).

Time-series irreversibility is measured by symbolic statistics \cite{daw2000}, by asymmetry indices of physiological series \cite{porta2008}, by visibility graphs \cite{lacasa2012}, also applied to financial series \cite{flanagan2016}, and by many other estimators \cite{zanin2021}; for linear processes reversibility criteria are classical \cite{weiss1975,tong2005,chan2006}, and the antisymmetric part of lagged cross-correlations is an established nonequilibrium measure, as the stochastic area or area-enclosing rate of linear systems \cite{ghanta2017,gonzalez2019} and as the cycling frequency of probability-flux analysis \cite{gradziuk2019,battle2016}. When only part of a system is observed, entropy production is bounded or inferred from coarse-grained trajectories \cite{roldan2010,roldan2012}, from waiting times and visible transitions \cite{skinner2021,harunari2022,vandermeer2022}, and exactly for coarse-grainings that preserve it \cite{teza2020}. These works ask how hidden degrees of freedom or projections change the measured irreversibility. Here the observable is a causal filter of the signal, and the question is how much irreversibility the temporal offsets of its kernels add.

In finance, time-reversal asymmetry of returns is a known stylized fact, traced to multi-timescale volatility feedback \cite{zumbach2009,lynch2003,chicheportiche2014} and to the leverage effect \cite{black1976,bouchaud2001}. What is new here is the exact split of the area-enclosing rate of causally filtered observables into an observer part, set by the kernel offsets, and a physical part, together with a model-free estimator from time-reversed data [Eq.~\eqref{eq:estimator}], its consequences for windowed indicators, and the measurement of both parts in an early-warning pipeline applied to market data. The split also shows which kind of financial time asymmetry a second-order trailing-window state can see: the leverage asymmetry between channels, not the higher-order asymmetry of volatility itself.
